\documentclass[11pt]{article}
\usepackage[margin=1in]{geometry}
\usepackage{amsmath,amssymb,amsthm}
\usepackage{booktabs}
\usepackage[T1]{fontenc}
\usepackage{lmodern}

\newtheorem{theorem}{Theorem}[section]
\newtheorem{lemma}[theorem]{Lemma}
\newtheorem{proposition}[theorem]{Proposition}
\newtheorem{corollary}[theorem]{Corollary}
\theoremstyle{definition}
\newtheorem{definition}[theorem]{Definition}
\theoremstyle{remark}
\newtheorem{remark}[theorem]{Remark}

\newcommand{\Ftwo}{\mathbb F_2}
\newcommand{\BH}{\mathcal B}
\newcommand{\EH}{\mathcal E}
\newcommand{\Cut}{\operatorname{Cut}}
\newcommand{\sfq}{\operatorname{sf}}
\newcommand{\supp}{\operatorname{supp}}
\newcommand{\rank}{\operatorname{rank}}

\title{The Signed-Cut Extinction Program for M\"obius\\
A New Finite-Combinatorial Route}
\author{V\'ictor Casa\~na--Giner\\
Bio Logic Crop Science S.L., E-46120 Alboraya (Valencia), Spain\\
\texttt{vcasana@biologiccropscience.com}\\
ORCID 0000-0002-8252-2608}
\date{}

\begin{document}
\maketitle

\begin{abstract}
We develop a finite combinatorial program for the M\"obius function based on the complete cube of squarefree multiplicative sign completions.  Each completion is encoded by a vector of odd-prime sign flips.  For every displacement $h$, the difference between its shifted correlation and the M\"obius correlation is exactly a signed cut in a displacement graph, and the admissible cuts are the image of a binary linear code generated by prime divisibility.  The edge signs themselves are determined by the Hamming parity of the associated prime-support form, so the survivor sets
\[
 \mathcal B_H(N)=\{t:d_h(t)\le0\text{ for }1\le h\le H\}
\]
are defined without using $C_h(\mu)$ as an external oracle.

The central finite phenomenon is complete extinction: exhaustive computations show that successive signed-cut inequalities eliminate every multiplicative completion except the M\"obius point $t=0$, often after very few lags.  This is stronger and conceptually different from eliminating only large-endpoint states.  The surviving M\"obius point is not discarded at the end of the argument: the extinction cells partition the entire nonzero prime-sign cube and yield an exact decomposition of the Mertens endpoint in terms of the eliminated states.  This is the continuation of the program.

A Fej\'er or moving-window mean-square estimate is recorded only as a non-route.  It can be obtained directly by applying the classical window inequality to $\mu$, without the completion cube, the signed cuts, or the extinction partition.  Therefore it bypasses the new mechanism and is not a continuation of the present proof program.  The route pursued here is the finite combinatorial one: prove complete extinction in general and extract the cancellation of $M(N)$ from the structure of the extinction partition itself.
\end{abstract}


\section{Introduction: a Darwinian view of the M\"obius word}

When confronting $M$, there is an immediate first intuition: being a fundamental, basic series, a driving force of all number theory, one must suspect (or even have the conviction) that $M$ is special.  I imagine the billions, trillions, infinitely many possible combinations with $1$, $-1$, and $0$, as competitors until reaching the mathematical ovum.  And the question is simple: what makes $M$ special so that it is the only one that survives in a selection process (in a Darwinian manner) until the end?  To answer this question there is an extremely simple process: to see which environmental and genetic conditions of $M$ (whose letters are not A, T, G, C but $1$, $-1$, and $0$) force its survival.  A great majority of methods, so far unsuccessful, focus on measuring the environmental conditions, and determine in what way they force the evolutionary genetic transformation of $M$.  But, as in genetics, that approach is very difficult (perhaps even impossible because of the lack of observables) to obtain or prove.  It would require knowing everything about the environmental conditions, which, like weather prediction, is unmanageable today.  However, if one takes the sequence of $M$ and modifies, one by one, a series of environmental factors, seeing how they affect $M$ and its competitors, and observes where its genetic code fails (almost in the manner of a triplet of a genetic codon) and why it dies, it is in principle possible to decipher not the genome of $M$ itself (which exists, is 100\% determined, and is theoretically computable to infinity), but why it has that genome.  And what makes impossible the, in the infinite, infinitely many combinations that never reach the ovum of mathematics.

The following work consists in determining that genetic code of $M$, its Darwinian evolution from species, subspecies, and nearby variants, with the conviction that such selection exists and is explainable.

The ingredients available to us are absolutely all of them: we have all the individuals of the species and we have absolutely exact sequencers, we know which lethal combinations are produced by certain genome recombinations, and all that remains is to determine, in a pool of Lego pieces in which we have all the pieces, which other combinations are lethal except one: the genome of $M$.

\subsection*{Mathematical formulation of the selection picture}

There is a first, almost unavoidable intuition when one confronts the Mertens sum and the M\"obius word that generates it: the M\"obius sequence should be treated as an exceptional object.  It is one of the most elementary arithmetic sequences and at the same time one of the basic motors of multiplicative number theory.  At a finite scale $N$, however, it sits inside an enormous population of possible words built from the alphabet
\[
 \{-1,0,+1\}.
\]
Once squarefree support and multiplicativity are imposed, the remaining population is still exponentially large: every choice of the odd-prime signs produces a competing multiplicative completion.  One may therefore picture billions, trillions, and ultimately an unbounded hierarchy of competitors undergoing a finite Darwinian selection process.  The elementary question is then not merely how large the Mertens sum is, but:
\[
 \boxed{\text{what makes the M\"obius word the unique survivor?}}
\]

The analogy is deliberately genetic.  The alphabet is not A, T, G, C but $-1,0,+1$; the heritable coordinates are the prime signs; recombination is multiplicative; and the successive displacement tests act as environmental pressures.  A competitor survives as long as every imposed test remains compatible with the M\"obius inequalities.  The first lag at which a positive signed cut appears is a lethal event for that competitor.  The resulting filtration is therefore a literal finite selection procedure: species, subspecies, and increasingly close variants are removed until only the M\"obius state remains.

A large part of the classical difficulty surrounding $M(N)$ comes from trying to infer the genetic outcome from the environment: zeros, primes, transforms, averages, and global analytic constraints are studied in order to deduce how the M\"obius sequence must behave.  That direction is extraordinarily demanding because it asks for sufficiently precise control of an environment with very few observables.  The present program reverses the direction.  Instead of attempting to model the full environment first, it changes the finite genetic coordinates one by one, or in controlled combinations, and records exactly where each competing genome becomes nonviable and why.

The point is not to reconstruct the M\"obius genome.  That genome already exists, is completely determined, and is algorithmically computable at every finite scale.  The problem is to explain \emph{why that genome is selected}.  Equivalently: among all multiplicative squarefree sign completions, which combinations are lethal, which survive a given collection of tests, and what structural obstruction prevents every non-M\"obius completion from reaching the final stage?

At finite $N$ all ingredients of this selection problem are available exactly.  The full population is explicit.  The prime-sign coordinates are exact.  The displacement graphs are exact.  The effect of each recombination is exact.  The signed-cut identity gives an exact sequencer for the first lethal defect.  Nothing statistical is required in the definition of the selection process.  One is therefore in the unusually favorable situation of having all the pieces of the combinatorial ``Lego pool'' and asking which assemblies are impossible except for one: the M\"obius word.

This paper develops that program.  Its central object is the nested family of survivor sets $\mathcal B_H(N)$.  The computational experiments motivating the theory do not merely remove half of the competitors or only the large-excursion states: they successively eliminate every admissible completion except the M\"obius point.  The purpose of the theory is to identify the exact finite combinatorial mechanism behind that extinction and then to use the extinction partition itself---rather than replace it by an unrelated classical criterion---to understand the endpoint $M(N)$ of the unique surviving genome.

\section{Scope: one route, and what is not the route}

The object of this paper is a deterministic finite problem.  At scale $N$ one considers all multiplicative squarefree sign completions of the M\"obius word, filters them by successive correlation inequalities, and studies the combinatorial information carried by the states eliminated at each lag.  The M\"obius word is the unique distinguished state $t=0$.

The continuation of this program is therefore internal:
\[
\boxed{
\begin{aligned}
\text{prime-sign cube}&\longrightarrow\text{ signed displacement cuts}\longrightarrow\mathcal B_H(N)\\
&\longrightarrow\text{ complete extinction of }t\ne0\\
&\longrightarrow\text{ endpoint information at }t=0.
\end{aligned}
}
\]
No step in this chain is to be replaced by a classical mean-square criterion for $M(N)$, by Fej\'er smoothing, or by a direct analytic reformulation of the Riemann hypothesis.  Such results may be used as comparisons, but not as substitutes for the combinatorial mechanism.

This restriction is mathematical rather than stylistic.  Section~\ref{sec:nonroute} proves that the Fej\'er/moving-window branch can be derived directly from $x=\mu$ without constructing $\mathcal B_H(N)$ at all.  Consequently, that branch discards exactly the finite extinction information introduced here.

\section{The finite multiplicative completion cube}

Let
\[
 Q(N):=\#\{n\le N:\mu(n)\ne0\}.
\]
For each odd prime $p\le N$, choose
\[
 \sigma_p\in\{-1,+1\},
\]
and fix
\[
 \sigma_2=-1.
\]
The associated admissible completion is
\[
 x_n=
 \begin{cases}
 \displaystyle\prod_{p\mid n}\sigma_p,&n\text{ squarefree},\\[1ex]
 0,&n\text{ not squarefree}.
 \end{cases}
\]
The M\"obius word is the state with $\sigma_p=-1$ for every prime.

Let
\[
 T:=\{p\le N:p\text{ odd prime and }\sigma_p=+1\},
\]
and let
\[
 t\in\Ftwo^{\pi(N)-1}
\]
be its indicator.  For squarefree $n$, let
\[
 v(n)\in\Ftwo^{\pi(N)-1}
\]
be the indicator vector of the odd prime factors of $n$.

\begin{lemma}[Admissible discrepancy sets are binary semispaces]\label{lem:semispace}
For every squarefree $n\le N$,
\[
 \frac{x_n}{\mu(n)}=(-1)^{\langle v(n),t\rangle}.
\]
Hence, if
\[
 D(t):=\{n\le N:x_n(t)\ne\mu(n)\},
\]
then
\[
 \boxed{
 n\in D(t)\iff \langle v(n),t\rangle=1.
 }
\]
\end{lemma}

\begin{proof}
On the squarefree support,
\[
 \mu(n)=\prod_{p\mid n}(-1),
\]
so
\[
 \frac{x_n}{\mu(n)}=\prod_{p\mid n}(-\sigma_p).
\]
The factor $-\sigma_p$ equals $-1$ exactly when $p\in T$.  The fixed prime $2$ contributes $+1$.  This gives the claimed parity formula.
\end{proof}

Thus the admissible discrepancy sets are not arbitrary subsets of the squarefree support.  They are precisely the codewords
\[
 \mathbf 1_{D(t)}=At\pmod2,
\]
where
\[
 A_{n,p}=\mathbf1_{p\mid n}\pmod2
\]
for odd primes $p$.

\section{Signed displacement graphs}

For $1\le h<N$, define
\[
 C_h(x):=\sum_{n\le N-h}x_nx_{n+h},
 \qquad
 d_h(t):=C_h(x_t)-C_h(\mu).
\]
Let $G_h(N)=(V_N,E_h,w_h)$, where
\[
 V_N:=\{n\le N:\mu(n)\ne0\},
\]
\[
 E_h:=\{\{n,n+h\}:1\le n\le N-h,\ n,n+h\in V_N\},
\]
and
\[
 w_h(\{n,n+h\}):=-\mu(n)\mu(n+h)\in\{\pm1\}.
\]
For an edge $e=\{n,n+h\}$ define its binary edge form
\[
 u_e:=v(n)\oplus v(n+h).
\]

\begin{theorem}[Exact signed-cut identity]\label{thm:cut}
For every admissible completion and every $1\le h<N$,
\[
 \boxed{
 d_h(t)=2\sum_{e\in\partial_hD(t)}w_h(e),
 }
\]
where $\partial_hD(t)$ is the set of edges having exactly one endpoint in $D(t)$.
Equivalently, if $A_h(t)$ and $B_h(t)$ are the numbers of crossing edges of weights $+1$ and $-1$, respectively, then
\[
 \boxed{d_h(t)=2(A_h(t)-B_h(t)).}
\]
\end{theorem}

\begin{proof}
Write $x_n=\varepsilon_n\mu(n)$ on $V_N$, with $\varepsilon_n=-1$ precisely on $D(t)$.  Then
\[
 x_nx_{n+h}-\mu(n)\mu(n+h)
 =\mu(n)\mu(n+h)(\varepsilon_n\varepsilon_{n+h}-1).
\]
The parenthesis is zero off the cut and $-2$ on the cut.  Summing and using $w_h=-\mu(n)\mu(n+h)$ proves the result.
\end{proof}

By Lemma~\ref{lem:semispace}, an edge crosses $D(t)$ exactly when
\[
 \langle u_e,t\rangle=1.
\]
Thus every lag is a weighted parity-threshold test on the prime-sign cube.

\subsection{The edge sign is determined by parity}

For $q\in\Ftwo^{\pi(N)-1}$, let
\[
 m_h(q):=\#\{e\in E_h:u_e=q\}.
\]

\begin{proposition}[Parity law for displacement edges]\label{prop:parity-law}
For every edge $e\in E_h$,
\[
 \boxed{
 w_h(e)=(-1)^{|u_e|+h+1}.
 }
\]
Consequently, for a fixed edge form $q$, all edges with $u_e=q$ have the same sign, and
\[
 \boxed{
 a_h(q):=\sum_{e:u_e=q}w_h(e)
 =m_h(q)(-1)^{|q|+h+1}.
 }
\]
\end{proposition}

\begin{proof}
Since both endpoints are squarefree,
\[
 \mu(n)\mu(n+h)=(-1)^{\omega(n)+\omega(n+h)}.
\]
Modulo $2$, every odd prime dividing both endpoints cancels in the sum of the two prime counts, leaving $|u_e|$.  The prime $2$ occurs in exactly one endpoint precisely when $h$ is odd.  Hence
\[
 \omega(n)+\omega(n+h)\equiv |u_e|+h\pmod2.
\]
Multiplication by the leading minus sign in the definition of $w_h$ gives the formula.
\end{proof}

This removes the apparent dependence on the M\"obius correlation values.  Indeed,
\[
 \sum_q a_h(q)=\sum_{e\in E_h}w_h(e)=-C_h(\mu),
\]
and therefore Theorem~\ref{thm:cut} becomes
\[
 \boxed{
 d_h(t)=\sum_q a_h(q)\bigl(1-(-1)^{\langle q,t\rangle}\bigr)
 =2\sum_{\langle q,t\rangle=1}a_h(q).
 }
 \tag{3.1}
\]
Combining this with Proposition~\ref{prop:parity-law},
\[
 \boxed{
 d_h(t)=
 2\sum_{\langle q,t\rangle=1}
 m_h(q)(-1)^{|q|+h+1}.
 }
 \tag{3.2}
\]
The right-hand side is determined by the finite factorization/divisibility geometry encoded in the multiplicities $m_h(q)$, the Hamming parities $|q|$, and the lag $h$.  It does not require $C_h(\mu)$ as an external input.

\begin{corollary}[First witness in oracle-free form]
For $t\ne0$, define
\[
 h_*(t):=\min\{h:d_h(t)>0\},
\]
when it exists.  Then
\[
 \boxed{
 h_*(t)=
 \min\left\{h:
 \sum_{\langle q,t\rangle=1}
 m_h(q)(-1)^{|q|+h+1}>0
 \right\}.
 }
\]
\end{corollary}

\subsection{The dead form and exact cube moments}

\begin{proposition}[The unique dead edge]\label{prop:dead}
The zero form occurs exactly as follows:
\[
 a_h(0)=
 \begin{cases}
 1,&h\text{ odd, squarefree, and }2h\le N,\\
 0,&\text{otherwise}.
 \end{cases}
\]
When it exists, the unique edge with $u_e=0$ is $\{h,2h\}$ and it has weight $+1$.  Since $\langle0,t\rangle=0$ for every $t$, this edge never crosses an admissible discrepancy cut.
\end{proposition}

\begin{proof}
The condition $u_e=0$ means that the two squarefree endpoints have the same odd-prime support.  Therefore they are $m$ and $2m$ for an odd squarefree $m$.  Their difference is $m$, so $m=h$ and the edge is unique.  Its weight is
\[
 -\mu(h)\mu(2h)=\mu(h)^2=1.
\]
\end{proof}

Let expectation below mean uniform averaging over the complete prime-sign cube.

\begin{proposition}[Exact first and second moments of the lag correlations]\label{prop:moments-c}
For every $h$,
\[
 \boxed{\mathbb E_t[C_h(x_t)]=-a_h(0),}
\]
and
\[
 \boxed{
 \operatorname{Var}_t[C_h(x_t)]
 =\sum_{q\ne0}a_h(q)^2.
 }
\]
Equivalently,
\[
 \operatorname{Var}_t[C_h(x_t)]
 =E_h-a_h(0)^2
 +\sum_{\substack{e\ne e'\\u_e=u_{e'}}}w_h(e)w_h(e'),
\]
where $E_h=|E_h|$.
\end{proposition}

\begin{proof}
From (3.1), or directly from the Walsh expansion,
\[
 C_h(x_t)=-\sum_q a_h(q)(-1)^{\langle q,t\rangle}.
\]
The characters of the binary cube are orthonormal.  Their mean is zero except for $q=0$, which proves the first identity; Parseval gives the variance.  Expanding $\sum_q a_h(q)^2$ by equal edge forms gives the final formula.
\end{proof}

When collisions $u_e=u_{e'}$ are sparse, the exact variance is of order $E_h$.  For fixed $h$, the squarefree-pair count satisfies $E_h\asymp N$, so the natural standard-deviation scale of $C_h$ over the completion cube is $\sqrt N$.

\subsection{Path structure}

\begin{proposition}[Each displacement graph is a forest of paths]\label{prop:paths}
Every connected component of $G_h(N)$ is a path.  In particular,
\[
 \boxed{
 \#\{\text{components of }G_h(N)\}=Q(N)-E_h.
 }
\]
\end{proposition}

\begin{proof}
Each vertex has at most the two possible neighbors $n-h$ and $n+h$.  A cycle cannot occur on the integer line with all edges of fixed positive length $h$.  Thus $G_h$ is a forest of paths.  A forest with $V$ vertices and $E$ edges has $V-E$ components.
\end{proof}

\section{The survivor filtration and complete extinction}

\begin{definition}[Survivors]
For $0\le H<N$, define
\[
 \boxed{
 \BH_H(N):=
 \{t\in\Ftwo^{\pi(N)-1}:d_h(t)\le0\text{ for every }1\le h\le H\}.
 }
\]
Then
\[
 \BH_0(N)\supseteq\BH_1(N)\supseteq\cdots\supseteq\BH_{N-1}(N).
\]
\end{definition}

The M\"obius state is $t=0$.  From (3.1),
\[
 d_h(0)=0
\]
for every lag, hence
\[
 \boxed{0\in\BH_H(N)\quad\text{for all }H<N.}
\]
The filtration can eliminate every competitor but can never eliminate M\"obius.

\begin{definition}[Complete extinction index]
When it exists, define
\[
 \boxed{
 H_0(N):=\min\{H<N:\BH_H(N)=\{0\}\}.
 }
\]
For $1\le h\le H_0(N)$ define the extinction cell
\[
 \boxed{
 \EH_h(N):=\BH_{h-1}(N)\setminus\BH_h(N).
 }
\]
\end{definition}

If $H_0(N)$ exists, then the entire non-M\"obius cube is partitioned by first positive witness:
\[
 \boxed{
 \Ftwo^{\pi(N)-1}\setminus\{0\}
 =\bigsqcup_{h=1}^{H_0(N)}\EH_h(N).
 }
 \tag{4.1}
\]
Equivalently,
\[
 t\in\EH_h(N)
 \iff
 d_j(t)\le0\ (j<h),\qquad d_h(t)>0.
\]
By (3.2), every cell in (4.1) is defined entirely by finite weighted parity thresholds.

\subsection{Exact finite benchmarks}

At $N=25$, the fixed-$\sigma_2$ family has $256$ states.  M\"obius is the unique state satisfying
\[
 C_1(x)\le C_1(\mu),\qquad
 C_2(x)\le C_2(\mu),\qquad
 C_3(x)\le C_3(\mu),
\]
so
\[
 H_0(25)=3.
\]

Further exhaustive computations give the following survivor counts.  Here $F=\pi(N)-1$ is the number of free odd-prime bits.

\begin{center}
\begin{tabular}{rrrrrrr}
\toprule
$N$ & $F$ & $|\BH_1|$ & $|\BH_2|$ & $|\BH_3|$ & $|\BH_5|$ & $H_0(N)$\\
\midrule
43 & 13 & 7345   & 2981   & 31   & 1   & 5\\
53 & 15 & 27384  & 8607   & 45   & 2   & 7\\
61 & 17 & 107624 & 35945  & 198  & 2   & 7\\
71 & 19 & 385588 & 58612  & 233  & 4   & 7\\
73 & 20 & 771176 & 162730 & 1691 & 97  & 9\\
\bottomrule
\end{tabular}
\end{center}

The phenomenon is not ``elimination of half the problem''.  At each tested scale the filtration removes the whole multiplicative completion cube except the single M\"obius state.  The central theorem sought by the program is the general finite-combinatorial analogue of this complete extinction.

\begin{theorem}[General complete extinction target]\label{thm:CE-target}
The combinatorial target is to prove, for the relevant range of $N$,
\[
 \boxed{
 \BH_{H_0(N)}(N)=\{0\}
 \quad\text{for some }H_0(N)<N,
 }
 \tag{CE}
\]
using only the weighted parity-threshold structure (3.2).
\end{theorem}

The statement is deliberately separated from any endpoint bound.  It says that the short-correlation order constraints uniquely determine the M\"obius point among all multiplicative squarefree sign completions.

\section{The endpoint is part of the same cube}

Define
\[
 S_N(t):=\sum_{n\le N}x_n(t).
\]
Since $\sigma_2=-1$, every odd squarefree $m\le N/2$ cancels with $2m$.  Therefore
\[
 \boxed{
 S_N(t)=
 \sum_{\substack{N/2<m\le N\\m\text{ odd, squarefree}}}
 (-1)^{\omega(m)}(-1)^{\langle v(m),t\rangle}.
 }
 \tag{5.1}
\]
In particular,
\[
 \boxed{S_N(0)=M(N).}
\]
Thus the Mertens endpoint is not an external analytic object attached after the extinction process.  It is the zero-character value of a Walsh polynomial on the very same prime-sign cube.

\begin{proposition}[Exact endpoint moments over the complete cube]\label{prop:endpoint-moments}
For $N\ge2$,
\[
 \boxed{\mathbb E_t[S_N(t)]=0,}
\]
and
\[
 \boxed{
 \operatorname{Var}_t[S_N(t)]
 =Q(N)-2Q_{\rm odd}(N/2),
 }
\]
where $Q_{\rm odd}(X)$ counts odd squarefree integers up to $X$.  Hence
\[
 \operatorname{Var}_t[S_N(t)]\sim\frac{2}{\pi^2}N.
\]
\end{proposition}

\begin{proof}
Every nonconstant character in (5.1) has mean zero, giving the first assertion.  Alternatively, in the unreduced sum the fixed terms $n=1$ and $n=2$ cancel and every term containing an odd prime has zero mean.

For the second moment, orthogonality gives
\[
 \mathbb E_t[S_N(t)^2]
 =\sum_{v(n)=v(m)}\mu(n)\mu(m).
\]
The diagonal contributes $Q(N)$.  The only distinct squarefree integers with the same odd-prime support are $k$ and $2k$ with $k$ odd squarefree.  Each unordered pair contributes
\[
 2\mu(k)\mu(2k)=-2.
\]
There are $Q_{\rm odd}(N/2)$ such pairs.  Since the mean is zero, the formula is the variance.  The asymptotic follows from
\[
 Q(N)\sim\frac6{\pi^2}N,
 \qquad
 Q_{\rm odd}(X)\sim\frac4{\pi^2}X.
\]
\end{proof}

The root-$N$ scale therefore already exists as an exact second-moment feature of the finite completion cube.  This does not by itself bound the distinguished state $t=0$; the purpose of complete extinction is to use the nontrivial geometry of the other states rather than replace it by an external mean-square theorem.

\section{The extinction decomposition of the Mertens endpoint}

The complete extinction partition provides an exact identity involving \emph{all} eliminated completions and the unique survivor.

\begin{theorem}[Extinction-cell decomposition of $M(N)$]\label{thm:cell-decomp}
Assume complete extinction (CE) at scale $N$.  Then
\[
 \boxed{
 M(N)
 =-\sum_{h=1}^{H_0(N)}\ \sum_{t\in\EH_h(N)}S_N(t).
 }
 \tag{6.1}
\]
Equivalently, if
\[
 \overline S_h(N):=\frac1{|\EH_h(N)|}\sum_{t\in\EH_h(N)}S_N(t),
\]
then
\[
 \boxed{
 M(N)=-\sum_{h=1}^{H_0(N)}|\EH_h(N)|\,\overline S_h(N).
 }
 \tag{6.2}
\]
\end{theorem}

\begin{proof}
Uniformly over the complete prime-sign cube,
\[
 \sum_{t\in\Ftwo^{\pi(N)-1}}S_N(t)=0
\]
by character orthogonality, or by Proposition~\ref{prop:endpoint-moments}.  Under complete extinction, the cube is the disjoint union of $\{0\}$ and the cells $\EH_h(N)$.  Since $S_N(0)=M(N)$,
\[
 0=M(N)+\sum_{h=1}^{H_0(N)}\sum_{t\in\EH_h(N)}S_N(t),
\]
which is (6.1).  Formula (6.2) is the same identity grouped by cell averages.
\end{proof}

This is the point at which the combinatorial program must continue.  The eliminated states are not disposable scaffolding.  Their first-witness cells carry the exact compensating endpoint whose total is $-M(N)$.

For each cell define the unnormalized Walsh coefficient
\[
 \widehat{\mathbf1_{\EH_h}}(q)
 :=\sum_{t\in\EH_h(N)}(-1)^{\langle q,t\rangle}.
\]
Substituting (5.1) into (6.1) yields the entirely finite identity
\[
 \boxed{
 M(N)=
 -\sum_{h=1}^{H_0(N)}
 \sum_{\substack{N/2<m\le N\\m\text{ odd, squarefree}}}
 (-1)^{\omega(m)}
 \widehat{\mathbf1_{\EH_h}}(v(m)).
 }
 \tag{6.3}
\]
Every set $\EH_h(N)$ in (6.3) is defined by the explicit threshold system (3.2).  Thus the remaining endpoint problem can be stated entirely inside the extinction geometry: determine the cancellation of the endpoint characters across the first-witness cells.

\begin{definition}[Internal combinatorial endpoint target]
The continuation of the program is to derive, from the threshold structure of the extinction cells themselves, a bound of RH scale for the right-hand side of (6.3): for every $\varepsilon>0$,
\[
 \boxed{
 \left|
 \sum_{h\le H_0(N)}
 \sum_{m\in U_N}
 (-1)^{\omega(m)}
 \widehat{\mathbf1_{\EH_h}}(v(m))
 \right|
 \ll_\varepsilon N^{1/2+\varepsilon},
 }
 \tag{CET}
\]
where $U_N$ is the odd squarefree upper-half set in (5.1).
\end{definition}

The target (CET) is not proposed as a new analytic equivalent to RH.  It is the exact endpoint of the same finite combinatorial elimination process: the cells are the states that the program itself removes, and the characters are the endpoint characters already present in the completion cube.

\section{Why the classical mean-square branch is not a continuation}\label{sec:nonroute}

There is a classical moving-window inequality.  For arbitrary $x\in\mathbb R^N$ and $1\le H<N$,
\[
 S(x)^2\le
 \frac{N+H}{H+1}
 \left[
 C_0(x)+2\sum_{h=1}^{H}
 \left(1-\frac{h}{H+1}\right)C_h(x)
 \right].
 \tag{7.1}
\]
It follows from Cauchy--Schwarz applied to all length-$(H+1)$ moving sums.  Applying (7.1) directly to $x=\mu$ gives
\[
 M(N)^2\le
 \frac{N+H}{H+1}\Phi_H(N),
 \tag{7.2}
\]
where
\[
 \Phi_H(N)
 =Q(N)+2\sum_{h=1}^{H}
 \left(1-\frac{h}{H+1}\right)C_h(\mu).
\]
By Fej\'er--Parseval,
\[
 \Phi_H(N)
 =\frac1{H+1}\sum_k
 \left(\sum_{k-H\le n\le k}\mu(n)\right)^2
\]
with the usual zero extension at the boundaries.

The key logical fact is immediate:

\begin{proposition}[The mean-square branch bypasses the extinction mechanism]
Inequality (7.2) and its Fej\'er--Parseval representation require neither the prime-sign completion cube, nor the signed-cut identity, nor the survivor filtration, nor complete extinction.  They follow by substituting $x=\mu$ directly into the classical window inequality.
\end{proposition}

Therefore a proof program that replaces the extinction cells by a theorem controlling $\Phi_H(N)$ has left the new combinatorial route.  It may be a valid classical route, but it is not a continuation of this one.  In particular, the long-studied mean-square problem for Mertens sums does not explain why the finite completion cube collapses to $t=0$, and proving such a mean-square estimate would not exploit the information generated by that collapse.

For this reason, the present program does \emph{not} list Fej\'er smoothing, terminal Mertens mean squares, or equivalent analytic RH criteria among its next steps.  The next step is (CE) together with the internal endpoint extraction (CET).

\section{Literature position and novelty}

The route introduced here is new.  The nearby literature contains several individual ingredients, but not the architecture developed above.

Aymone and Sidoravicius study random multiplicative functions supported on squarefree integers with independent prime signs and analyze their partial sums.  That is a probability model, not the deterministic finite cube filtered by successive correlation inequalities.  Koukoulopoulos and Soundararajan, and Klurman--Mangerel--Pohoata--Ter\"av\"ainen, study structural and inverse questions for multiplicative functions with small partial sums, but not finite M\"obius completions, first-witness cells, or simultaneous signed displacement cuts.  Bourgain studies the Fourier--Walsh spectrum of $\mu$ on the Boolean cube of binary digits of the integer argument; this is a different Boolean cube from the prime-sign completion cube used here.  Matom\"aki--Radziwi\l\l--Tao study sign patterns and shifted behavior of M\"obius and Liouville directly, not an adversarial family of multiplicative completions filtered to a unique survivor.

None of these lines implements the defining chain
\[
\boxed{
\begin{aligned}
&\text{finite prime-sign completion cube}
\to\text{ signed displacement cuts}
\to\BH_H(N)\\
&\qquad\to\text{ first-witness extinction cells}
\to\text{ unique survivor }t=0\\
&\qquad\to\text{ endpoint extraction from the eliminated cells}.
\end{aligned}
}
\]
The novelty of the program lies in this chain.  Generic signed graphs, Walsh characters, shifted correlations, and Fej\'er identities are classical objects individually; their existence does not make the extinction architecture classical.

The route is also structurally distinct from the classical analytic approaches to RH.  The latter begin with zeta, Dirichlet series, zero-free regions, moments, or direct estimates for $M(x)$.  The present route begins with a finite constraint-satisfaction problem on prime signs and seeks to recover $M(N)$ only after all competing multiplicative completions have been combinatorially eliminated.

\section{What is proved, what is computational, and what remains}

The following statements are proved exactly in this paper:
\begin{itemize}
\item admissible discrepancy sets are binary semispaces (Lemma~\ref{lem:semispace});
\item correlation differences are signed displacement cuts (Theorem~\ref{thm:cut});
\item edge signs are determined by $|q|$ and $h$, giving the oracle-free threshold formula (3.2) (Proposition~\ref{prop:parity-law});
\item the zero edge form is the unique dead edge $\{h,2h\}$ when it exists (Proposition~\ref{prop:dead});
\item the first and second lag-correlation moments over the full completion cube are exact (Proposition~\ref{prop:moments-c});
\item each $G_h$ is a forest of paths (Proposition~\ref{prop:paths});
\item the endpoint $S_N(t)$ is a Walsh polynomial on the same cube and has exact variance $Q(N)-2Q_{\rm odd}(N/2)$ (Proposition~\ref{prop:endpoint-moments});
\item whenever complete extinction holds, the eliminated first-witness cells give the exact Mertens decomposition (6.1)--(6.3) (Theorem~\ref{thm:cell-decomp});
\item the classical Fej\'er/moving-window branch bypasses the extinction mechanism and therefore is not a continuation of the present route.
\end{itemize}

The complete collapse
\[
 \BH_{H_0(N)}(N)=\{0\}
\]
is verified exhaustively at the finite scales displayed above, including $N=25,43,53,61,71,73$.  At those scales every multiplicative completion except M\"obius is eliminated by a positive signed-cut witness.  This is stronger than elimination of an RHC or large-endpoint subfamily.

The remaining mathematical problem is wholly combinatorial: prove complete extinction in general and exploit the explicit first-witness cells to control the unique surviving endpoint through (6.3).  No classical mean-square theorem is designated as the next step, because such a theorem can be formulated without the completion cube and would abandon the information generated by extinction.

\section{Conclusion}

The finite completion experiment does not remove only a selected half of the admissible states.  It removes every tested non-M\"obius completion.  The corresponding exact formalism is a simultaneous signed-cut filtration on a prime-sign Boolean cube, with no external correlation oracle after the parity reduction.  When complete extinction occurs, all nonzero states are partitioned into first-witness cells, and those same cells provide an exact decomposition of $M(N)$.

This identifies a single continuation of the program: prove the extinction geometry and use it all the way through the endpoint.  Replacing that step by Fej\'er--Parseval or a classical Mertens mean-square estimate would discard the new combinatorial information rather than complete it.  The route developed here is therefore a distinct new finite-combinatorial route to the Mertens/RH problem, and its unresolved step is internal to that route.

\begin{thebibliography}{99}

\bibitem{AymoneSidoravicius}
M. Aymone and V. Sidoravicius,
\emph{Partial sums of biased random multiplicative functions},
J. Number Theory 172 (2017), 343--382;
arXiv:1412.1157.

\bibitem{AymoneBuenoMedeiros}
M. Aymone, C. Bueno, and K. Medeiros,
\emph{Multiplicative functions supported on the $k$-free integers with small partial sums},
arXiv:2101.00279.

\bibitem{Bourgain}
J. Bourgain,
\emph{On the Fourier--Walsh spectrum of the M\"obius function},
arXiv:1112.1423.

\bibitem{KoukSound}
D. Koukoulopoulos and K. Soundararajan,
\emph{The structure of multiplicative functions with small partial sums},
arXiv:1909.13105.

\bibitem{KMPT}
O. Klurman, A. P. Mangerel, C. Pohoata, and J. Ter\"av\"ainen,
\emph{Multiplicative functions that are close to their mean},
arXiv:1911.06265.

\bibitem{MRT}
K. Matom\"aki, M. Radziwi\l\l, and T. Tao,
\emph{Sign patterns of the Liouville and M\"obius functions},
arXiv:1509.01545.

\end{thebibliography}

\end{document}
